\sloppy
\emergencystretch=10em
\section{Overview}

This chapter presents a post-training, weight-only quantization methodology for the query, key, and value projections of transformer attention. The methodology is designed for asymmetric group-wise INT4 quantization: projection weights are stored on a four-bit integer grid with group-specific scales and zero points, whereas activations remain in their original precision \cite{jacob2018quantization,nagel2021white}. The central premise is that nearest rounding is a strong starting point but not necessarily the best final choice when the quantized weights are evaluated through the activations and interactions that determine attention.

The proposed refinement has two complementary stages. The first stage, called \textbf{Flip}, treats a projection matrix independently and reduces its activation-weighted output residual. Flip is deliberately distinguished from ordinary round-to-nearest (RTN) quantization. RTN already minimizes the isolated scalar reconstruction error of each weight on a fixed quantization grid. Consequently, changing a rounded integer to a neighboring level generally cannot improve that isolated scalar objective. Flip instead accepts a neighboring level only when the resulting change reduces the aggregate projected residual. Such a change may increase the error of an individual weight while improving the output of the complete projection through cancellation among activation-weighted errors.

The second stage, called \textbf{ReFlip}, addresses the coupling between query and key projections within grouped-query attention (GQA). For each query head, ReFlip holds its group's shared quantized key head fixed and updates only that query head's quantized query matrix. It evaluates how a legal one-step integer change affects an unscaled representative-vector query--key surrogate. The resulting objective is exact for this scalar surrogate: no Taylor approximation or full-matrix gradient surrogate is used. ReFlip can therefore compensate for a surrogate residual jointly caused by query and key quantization, even though only query weights are modified. This compensation is useful, but its scope must be stated precisely: preservation of one representative scalar surrogate does not imply preservation of the full pairwise attention-score matrix over all token pairs.

Both refinement stages use a representative activation row vector obtained from calibration data and stabilized by James--Stein shrinkage. This produces a compact, data-dependent signal for evaluating discrete changes without retraining. Kneedle-based selection is used in ReFlip to identify an extreme high-sensitivity region that should be protected and a moderate-sensitivity region immediately after the knee in which correction is permitted. Kneedle is an adaptive heuristic rather than an optimality guarantee.

The overall procedure remains fully post-training. It requires neither gradient-based optimization nor modification of the model architecture. The refinement is performed before deployment, and the resulting matrices retain the same group-wise asymmetric INT4 representation used by the initial quantizer. Thus, the method aims to improve the functional fidelity of quantized attention projections while preserving the principal memory motivation of low-bit weight storage and without making claims about measured runtime.

\section{Notation and Calibration Representative}

\subsection{Projection Convention}

Let a representative activation be a row vector
\begin{equation}
x \in \mathbb{R}^{d_{\mathrm{in}}}.
\end{equation}
For a generic projection, let the full-precision weight matrix have shape
\begin{equation}
W \in \mathbb{R}^{d_{\mathrm{out}} \times d_{\mathrm{in}}},
\end{equation}
so that the projected row vector is
\begin{equation}
y=xW^\top\in\mathbb{R}^{d_{\mathrm{out}}}.
\end{equation}
This row-vector convention is used throughout the chapter. In particular, matrix row $a$ identifies an output dimension and matrix column $b$ identifies an input dimension.

The quantized and dequantized approximation is denoted by
\begin{equation}
\hat W\in\mathbb{R}^{d_{\mathrm{out}} \times d_{\mathrm{in}}},
\end{equation}
with projected vector
\begin{equation}
\hat y=x\hat W^\top.
\end{equation}
Although the stored representation is integer-valued, the hat notation refers to the dequantized value represented by those integers, scales, and zero points.

For GQA \cite{ainslie2023gqa}, let $\mathcal{H}_Q$ be the set of query-head indices and $\mathcal{H}_{KV}$ the set of key/value-head indices. The query heads are partitioned into GQA groups
\begin{equation}
\{\mathcal{G}_r:r\in\mathcal{H}_{KV}\},
\qquad
\bigcup_{r\in\mathcal{H}_{KV}}\mathcal{G}_r=\mathcal{H}_Q,
\qquad
\mathcal{G}_r\cap\mathcal{G}_{r'}=\varnothing\quad(r\neq r'),
\end{equation}
where every query head $h\in\mathcal{G}_r$ shares key/value head $r$. With common head dimension $d_h$, the head-specific projection matrices are
\begin{equation}
W_Q^{(h)},\,W_K^{(r)},\,W_V^{(r)}
\in\mathbb{R}^{d_h\times d_{\mathrm{in}}},
\qquad h\in\mathcal{G}_r.
\end{equation}
Their representative projected vectors are
\begin{equation}
q^{(h)}=x(W_Q^{(h)})^\top,\qquad
k^{(r)}=x(W_K^{(r)})^\top,\qquad
v^{(r)}=x(W_V^{(r)})^\top,
\end{equation}
and the corresponding dequantized vectors are written $\hat q^{(h)}$, $\hat k^{(r)}$, and $\hat v^{(r)}$. This notation distinguishes a GQA group, which describes head sharing, from a quantization group, which describes columns sharing quantization parameters.

\subsection{James--Stein-Shrunk Activation}

Calibration activations can be noisy when only a limited post-training sample is available. A single unregularized statistic may overemphasize accidental extremes, whereas a simple global average may suppress meaningful differences between input dimensions. To obtain a regularized representative row vector, this work uses a nonnegative James--Stein-inspired shrinkage heuristic, while retaining the distinction between this calibration construction and the assumptions of the classical estimator \cite{james1961estimation}.

Let $X^{(1)},\ldots,X^{(N_{\mathrm{cal}})}\in\mathbb{R}^{d_{\mathrm{in}}}$ denote the calibration activation rows. The empirical channel-mean vector and its grand mean are
\begin{equation}
\tilde{x}_b=\frac{1}{N_{\mathrm{cal}}}\sum_{n=1}^{N_{\mathrm{cal}}}X_b^{(n)},
\qquad
\bar{x}=\frac{1}{d_{\mathrm{in}}}\sum_{b=1}^{d_{\mathrm{in}}}\tilde{x}_b.
\end{equation}
The shrinkage target is the constant vector $m=\bar{x}\mathbf{1}$, and the calibration-noise proxy is the squared mean absolute deviation of the channel means,
\begin{equation}
\hat{\sigma}^2=
\left(
\frac{1}{d_{\mathrm{in}}}
\sum_{b=1}^{d_{\mathrm{in}}}
\left|\tilde{x}_b-\bar{x}\right|
\right)^2.
\end{equation}
The representative activation is
\begin{equation}
x=m+\lambda(\tilde{x}-m),
\end{equation}
where the positive-part shrinkage factor has the form
\begin{equation}
\lambda=
\left[
1-\frac{(d_{\mathrm{in}}-2)\hat{\sigma}^{2}}
{\|\tilde{x}-m\|_2^2}
\right]_+,
\qquad [u]_+=\max(u,0).
\end{equation}
If $\|\tilde{x}-m\|_2^2=0$, the representative is defined directly as $x=\tilde{x}=m$ and no shrinkage calculation is required. Otherwise, if the observed dimension-wise variation is weak relative to the noise proxy, $\lambda$ moves the statistic toward $m$; if the observed variation is strong, more of the empirical structure is retained. Because both $m$ and $\hat{\sigma}^2$ are estimated from the same calibration observations and the noise proxy is heuristic, this construction is inspired by the classical estimator rather than an invocation of its classical risk-dominance theorem. The resulting $x$ is not claimed to represent every token or every sequence. It is a compact calibration summary used to define the exact objectives evaluated by Flip and ReFlip.

This distinction is important. The refinement stages target behavior under the chosen representative activation. Their objectives can improve a targeted calibration criterion, but they do not provide a universal guarantee over the complete activation distribution.

\section{Problem Formulation}

The goal is to construct group-wise asymmetric INT4 approximations of the attention projection matrices while limiting functionally relevant quantization error. Conventional RTN quantization primarily considers each scalar weight on its own. Attention, however, is determined by projected activations and by interactions between the resulting query and key vectors.

For a generic projection matrix $W$, define the dequantization error
\begin{equation}
E=\hat W-W.
\end{equation}
Under the representative activation, the corresponding projection residual is
\begin{equation}
r=xE^\top=x(\hat W-W)^\top
\in\mathbb{R}^{d_{\mathrm{out}}}.
\end{equation}
The residual in output dimension $j$ is therefore
\begin{equation}
r_j=\sum_{i=1}^{d_{\mathrm{in}}}x_iE_{ji}.
\end{equation}
This expression makes the relevant error aggregation explicit. Individual signed weight errors are multiplied by activation coordinates and then summed. Two nonzero scalar weight errors can partially cancel in the projected output, while a small scalar error can matter substantially when multiplied by a large activation coordinate.

The local functional objective used by Flip is
\begin{equation}
J(W,\hat W;x)=\|r\|_2^2
=\left\|x(\hat W-W)^\top\right\|_2^2.
\end{equation}
This objective is applicable independently to $W_Q$, $W_K$, and $W_V$. It measures projection fidelity for the representative activation and does not claim to optimize the full attention computation.

The interaction-aware objective used by ReFlip is narrower and more explicit. For query head $h\in\mathcal{G}_r$, define the original unscaled representative-vector query--key surrogate
\begin{equation}
s_{h,r}^\star
=
\left[x(W_Q^{(h)})^\top\right]\cdot
\left[x(W_K^{(r)})^\top\right]
=q^{(h)}\cdot k^{(r)},
\end{equation}
and the current quantized surrogate
\begin{equation}
\hat{s}_{h,r}
=
\left[x(\hat W_Q^{(h)})^\top\right]\cdot
\left[x(\hat W_K^{(r)})^\top\right]
=\hat q^{(h)}\cdot\hat k^{(r)}.
\end{equation}
ReFlip uses the per-query-head loss
\begin{equation}
\mathcal{L}_{h,r}
=\frac{1}{2}(\hat{s}_{h,r}-s_{h,r}^\star)^2.
\end{equation}
The scalar residual
\begin{equation}
e_{h,r}=\hat{s}_{h,r}-s_{h,r}^\star
\end{equation}
contains the combined effect of quantizing query head $h$ and its shared key head $r$. ReFlip holds $\hat W_K^{(r)}$ fixed and changes only $\hat W_Q^{(h)}$, using the latter as the available degree of freedom for reducing this joint surrogate residual. The word ``unscaled'' is deliberate: this compact surrogate omits the usual $1/\sqrt{d_h}$ factor because multiplication by a positive constant would not change which exact one-step candidate decreases the squared residual. It is not presented as a general or complete attention logit.

These two objectives establish a progression from projection-level fidelity to query--key interaction fidelity. Flip targets an activation-weighted residual for one matrix at a time. ReFlip then compensates the remaining per-query-head scalar Q--K surrogate residual through legal integer changes to $\hat W_Q^{(h)}$. Neither objective is presented as an exact surrogate for all downstream model behavior.

\section{Group-Wise Asymmetric INT4 Quantization}

\subsection{Integer Representation}

For a generic projection matrix
\begin{equation}
W\in\mathbb{R}^{d_{\mathrm{out}}\times d_{\mathrm{in}}},
\end{equation}
the input dimension is partitioned into contiguous quantization groups of at most $G$ columns. Let $\gamma(b)$ map input column $b$ to its quantization-group index. Each output row and quantization group has a scale
\begin{equation}
\alpha_{a,\gamma(b)}>0
\end{equation}
and a zero point
\begin{equation}
\zeta_{a,\gamma(b)}\in\mathcal{Z}_4.
\end{equation}
The stored INT4 value at row $a$ and column $b$ is
\begin{equation}
z_{ab}\in\mathcal{Z}_4=\{0,1,\ldots,15\}.
\end{equation}
Its dequantized value is
\begin{equation}
\hat W_{ab}
=\alpha_{a,\gamma(b)}
\left(z_{ab}-\zeta_{a,\gamma(b)}\right).
\end{equation}

For row $a$ and quantization group $t$, let $\mathcal{I}_t=\{b:\gamma(b)=t\}$ be its input-column set. The real range is enlarged to include zero:
\begin{equation}
w^-_{a,t}
=
\min\left(0,\min_{b\in\mathcal{I}_t}W_{ab}\right),
\qquad
w^+_{a,t}
=
\max\left(0,\max_{b\in\mathcal{I}_t}W_{ab}\right).
\end{equation}
When the augmented range is nondegenerate, the asymmetric parameters are
\begin{equation}
\alpha_{a,t}
=
\frac{w^+_{a,t}-w^-_{a,t}}{15},
\qquad
\zeta_{a,t}
=
\operatorname{clip}_{[0,15]}
\left(
\operatorname{round}
\left(
-\frac{w^-_{a,t}}{\alpha_{a,t}}
\right)
\right).
\end{equation}
Including zero handles one-sided groups explicitly. If all weights in a group are nonnegative, then $w^-_{a,t}=0$ and $\zeta_{a,t}=0$; if all are nonpositive, then $w^+_{a,t}=0$ and the zero point is placed at the upper end of the integer range, up to the stated rounding and clipping rule. For a mixed-sign group, the zero point lies inside the legal range.

The initial RTN integer is
\begin{equation}
z_{ab}^{\mathrm{RTN}}
=\operatorname{clip}_{[0,15]}
\left(
\operatorname{round}
\left(
\frac{W_{ab}}{\alpha_{a,\gamma(b)}}+
\zeta_{a,\gamma(b)}
\right)
\right).
\end{equation}
Under the zero-inclusive range, a degenerate group satisfies $w^-_{a,t}=w^+_{a,t}=0$ and therefore contains only zeros. Such a group is represented exactly by the explicit convention
\begin{equation}
\alpha_{a,t}=1,\qquad
\zeta_{a,t}=0,\qquad
z_{ab}=0\quad\text{for all }b\in\mathcal{I}_t,
\end{equation}
and is frozen during refinement because every useful target value in the group is already represented exactly. This convention avoids division by zero without assigning an arbitrary effective step to an active nonzero group.

\subsection{What RTN Optimizes}

For fixed $\alpha_{a,\gamma(b)}$ and $\zeta_{a,\gamma(b)}$, RTN chooses the legal quantization level closest to $W_{ab}$. Including legal endpoint levels, it solves the isolated scalar problem, with non-uniqueness only at exact ties:
\begin{equation}
\min_{z\in\mathcal{Z}_4}
\left|
\alpha_{a,\gamma(b)}
\left(z-\zeta_{a,\gamma(b)}\right)-W_{ab}
\right|.
\end{equation}
Equivalently, it minimizes the squared scalar reconstruction error on the fixed grid. Therefore, a neighboring integer change should not be described as improving the isolated scalar weight error that RTN already minimizes.

The limitation of RTN is instead an objective mismatch. The projected residual in output row $a$ depends on the signed sum
\begin{equation}
r_a=\sum_b x_b(\hat W_{ab}-W_{ab}),
\end{equation}
not on each absolute scalar error separately. A neighboring integer can increase
\begin{equation}
|\hat W_{ab}-W_{ab}|
\end{equation}
yet reduce $|r_a|$ because its signed activation-weighted contribution better cancels the contributions of other weights. Flip exploits precisely this possibility.

\subsection{Memory and Inference Motivation}

INT4 storage provides the theoretical basis for reducing the weight memory of QKV projections. Ignoring metadata, four-bit weights require one quarter of the storage of FP16 weights and one eighth of the storage of FP32 weights. Group-wise scales and zero points introduce overhead, so the exact compression ratio depends on group size and metadata precision. Smaller groups can improve local representational fidelity but increase this overhead.

The refinement stages do not change the deployed tensor shapes or the mathematical dequantization rule. They modify legal stored integers while retaining the same scales and zero points. Consequently, the refined model remains compatible with the same class of group-wise weight-only inference procedure as its RTN initialization. Because the present methodology does not report a measured runtime study, it makes no empirical latency or throughput claim. Its inference motivation is structural: low-bit weight storage can reduce memory traffic and capacity requirements, while the offline refinement seeks to recover fidelity without introducing a new runtime objective or a new projection operator.

\section{Flip: Activation-Weighted Projection Refinement}

\subsection{Flip Objective}

For one projection matrix $W$ and its current quantized approximation $\hat W$, Flip targets
\begin{equation}
J=\|r\|_2^2,
\qquad
r=xE^\top,
\qquad
E=\hat W-W.
\end{equation}
The objective separates across output rows:
\begin{equation}
J=\sum_{j=1}^{d_{\mathrm{out}}}r_j^2.
\end{equation}
Nevertheless, it does not separate across the weights within a row because all input coordinates contribute to the same scalar residual $r_j$. This within-row coupling is the source of useful error cancellation.

Suppose a legal candidate changes only the dequantized weight at row $j$ and column $i$ by
\begin{equation}
\delta=\hat W'_{ji}-\hat W_{ji}.
\end{equation}
For a one-step integer move,
\begin{equation}
\delta=\alpha_{j,\gamma(i)}\Delta z,
\qquad
\Delta z\in\{-1,+1\},
\end{equation}
provided that $z_{ji}+\Delta z\in\mathcal{Z}_4$. The candidate changes only residual component $r_j$:
\begin{equation}
r'_j=r_j+x_i\delta.
\end{equation}
All other residual components remain unchanged. The exact objective change is consequently
\begin{align}
\Delta J
&=J'-J \\
&=(r_j+x_i\delta)^2-r_j^2 \\
&=2r_jx_i\delta+(x_i\delta)^2.
\end{align}
No approximation is required. A legal candidate is accepted only if
\begin{equation}
\Delta J<0.
\end{equation}
After acceptance, the scalar residual for that output row is updated as
\begin{equation}
r_j\leftarrow r_j+x_i\delta.
\end{equation}
This incremental update preserves exact agreement with the current quantized matrix while avoiding repeated reconstruction of the full residual.

\subsection{Interpretation as Error Cancellation}

The sign of the first term in $\Delta J$ determines whether a candidate initially moves against the current residual:
\begin{equation}
2r_jx_i\delta<0.
\end{equation}
The second term,
\begin{equation}
(x_i\delta)^2,
\end{equation}
penalizes an excessively large move and prevents a candidate from being accepted merely because it has the correct direction. The exact acceptance condition can also be written as
\begin{equation}
|r_j+x_i\delta|<|r_j|.
\end{equation}
Thus, Flip selects a neighboring quantization level only when its activation-weighted contribution reduces the magnitude of the aggregate row residual.

This interpretation resolves an apparent contradiction between RTN and Flip. RTN is already optimal for the isolated scalar reconstruction error on the fixed grid. Flip does not overturn that fact. Instead, it changes the objective from isolated scalar error to activation-weighted projection error. Consider a row containing two signed weight errors whose activation-weighted contributions reinforce one another. Moving one weight away from its nearest level can reverse or weaken one contribution, producing greater cancellation in their sum. The modified weight can be less accurate in isolation while the projected output becomes more accurate.

Accordingly, Flip must not be characterized as a second nearest-rounding search. It is a discrete coordinate refinement under the exact aggregate objective $J$. Its benefit depends on the representative activation and on the signed configuration of errors within each output row.

\subsection{Candidate Evaluation and Acceptance}

Flip begins from the RTN-quantized matrix and its exact residual $r$. For each eligible integer $z_{ji}$, it evaluates the legal directions in
\begin{equation}
\mathcal{D}_{ji}
=\left\{
\Delta z\in\{-1,+1\}:
z_{ji}+\Delta z\in\mathcal{Z}_4
\right\}.
\end{equation}
Here an eligible coordinate is one whose quantization group is nondegenerate and not frozen, and whose legal direction set $\mathcal{D}_{ji}$ is nonempty. Frozen degenerate groups and coordinates with no legal one-level move are skipped rather than assigned an artificial step size.
For each direction, it obtains $\delta=\alpha_{j,\gamma(i)}\Delta z$ and evaluates the exact $\Delta J$. If neither direction decreases $J$, the integer remains unchanged. When both symmetric directions are legal and $x_i\alpha_{j,\gamma(i)}\neq0$, at most one can strictly reduce $J$: adding their two exact objective changes yields the positive quantity $2(x_i\alpha_{j,\gamma(i)})^2$. Thus, the method accepts the unique strictly improving direction when one exists. Once a change is accepted, both $z_{ji}$ and $r_j$ are updated before subsequent candidates are evaluated.

Updating the residual after every accepted change is essential. Candidate gains are not independent: a previous change alters the cancellation needed in row $j$, and therefore alters the exact value of every later $\Delta J$ in that row. Evaluating all candidates against a stale initial residual would not implement the stated greedy objective.

Candidate order is a heuristic choice. One option is to evaluate coordinates in descending order of the magnitude of their possible activation-weighted step,
\begin{equation}
|x_i|\alpha_{j,\gamma(i)}.
\end{equation}
Another is to scan deterministically by row and column. The acceptance equation remains exact under either ordering, but the final discrete configuration can differ because accepted changes modify later residuals. A pass budget and an acceptance budget are therefore explicit hyperparameters rather than hidden assumptions.

\subsection{Scope and Limitations of Flip}

Flip can be applied independently to the query, key, and value projection matrices because its objective requires only a matrix, its quantized approximation, and the representative activation. For $W_Q$ and $W_K$, reducing the projection residual can indirectly improve the vectors used in query--key scoring. For $W_V$, it can improve the representative value projection. However, Flip does not directly optimize the interaction between query and key, the softmax operation, or the final attention output.

The objective is also calibration-dependent. A change that reduces $J$ under the representative James--Stein-shrunk vector may not reduce projection error for every activation row. The shrinkage estimate and conservative integer neighborhood are intended to limit sensitivity to calibration noise, but they do not constitute a distribution-wide guarantee. Finally, the greedy coordinate procedure can stop at a discrete local configuration even when a coordinated multi-weight change would yield a lower objective.

\section{ReFlip: Scalar Query--Key Surrogate Refinement}

\subsection{Per-Query-Head Fixed-Key Objective}

After initial quantization and the Flip stage in the proposed refinement sequence, ReFlip treats each query head $h\in\mathcal{G}_r$ separately while using the fixed shared key head $r$. The full-precision and current quantized unscaled representative-vector Q--K surrogates are
\begin{equation}
s_{h,r}^\star
=q^{(h)}\cdot k^{(r)}
\end{equation}
and
\begin{equation}
\hat{s}_{h,r}
=\hat q^{(h)}\cdot\hat k^{(r)}.
\end{equation}
Define the per-head scalar residual and loss
\begin{equation}
e_{h,r}=\hat{s}_{h,r}-s_{h,r}^\star,
\qquad
\mathcal{L}_{h,r}=\frac{1}{2}e_{h,r}^2.
\end{equation}

During the correction of query head $h$, $\hat W_K^{(r)}$ is fixed. ReFlip updates $\hat W_Q^{(h)}$ only; it does not update any key or value matrix. Although the available changes belong only to one query head, $e_{h,r}$ includes quantization effects from both that query head and its shared key head. ReFlip may therefore use a query-weight change to compensate part of the joint Q--K surrogate residual.

This scalar is explicitly an unscaled representative-vector surrogate. It is narrower than a transformer's complete attention computation: it uses one calibration representative on both projection paths, omits cross-token query--key pairs, and does not include softmax or value aggregation. Consequently, reducing $\mathcal{L}_{h,r}$ does not assert preservation of the full attention-score matrix or of a general attention logit.

\subsection{Exact Effect of an Integer Change}

Consider integer $z^{(h)}_{ab}$ representing query-head weight $(a,b)$. ReFlip directly changes
\begin{equation}
z^{(h)}_{ab}\leftarrow z^{(h)}_{ab}+\Delta z^{(h)}_{ab},
\qquad
\Delta z^{(h)}_{ab}\in\{-1,+1\},
\end{equation}
subject to
\begin{equation}
z^{(h)}_{ab}+\Delta z^{(h)}_{ab}\in\mathcal{Z}_4.
\end{equation}
Because the query-head scale and zero point remain fixed, the dequantized weight change is
\begin{equation}
\Delta\hat W^{(h)}_{Q,ab}
=\alpha^{(h)}_{a,\gamma(b)}\Delta z^{(h)}_{ab}.
\end{equation}
Only query coordinate $\hat q^{(h)}_a$ changes:
\begin{equation}
\Delta\hat q^{(h)}_a
=x_b\alpha^{(h)}_{a,\gamma(b)}\Delta z^{(h)}_{ab}.
\end{equation}
Since $\hat k^{(r)}$ is fixed, the exact surrogate change is
\begin{equation}
\Delta s^{(h,r)}_{ab}
=\hat k^{(r)}_a x_b
\alpha^{(h)}_{a,\gamma(b)}
\Delta z^{(h)}_{ab}.
\end{equation}

This expression is exact for the stated scalar surrogate, rather than a first-order approximation. Its exact loss change is
\begin{align}
\Delta\mathcal{L}^{(h,r)}_{ab}
&=
\frac{1}{2}
\left(e_{h,r}+\Delta s^{(h,r)}_{ab}\right)^2
-\frac{1}{2}e_{h,r}^2 \\
&=
e_{h,r}\Delta s^{(h,r)}_{ab}
+\frac{1}{2}\left(\Delta s^{(h,r)}_{ab}\right)^2.
\end{align}
A candidate is accepted only if
\begin{equation}
\Delta\mathcal{L}^{(h,r)}_{ab}<0.
\end{equation}
After acceptance, the integer parameter, dequantized query weight, query coordinate, scalar surrogate, and residual are updated immediately:
\begin{equation}
z^{(h)}_{ab}
\leftarrow
z^{(h)}_{ab}+\Delta z^{(h)}_{ab},
\qquad
\hat W^{(h)}_{Q,ab}
\leftarrow
\hat W^{(h)}_{Q,ab}+\Delta\hat W^{(h)}_{Q,ab},
\end{equation}
\begin{equation}
\hat q^{(h)}_a
\leftarrow
\hat q^{(h)}_a+\Delta\hat q^{(h)}_a,
\qquad
\hat{s}_{h,r}
\leftarrow
\hat{s}_{h,r}+\Delta s^{(h,r)}_{ab},
\end{equation}
\begin{equation}
e_{h,r}
\leftarrow
e_{h,r}+\Delta s^{(h,r)}_{ab}.
\end{equation}
The immediate residual update is necessary because every subsequent exact loss change is evaluated against the current quantized state.

\subsection{Integer-Level and Dimension Sensitivity}

The magnitude of surrogate change per unit change in the integer parameter directly modified by ReFlip defines the integer-level sensitivity:
\begin{equation}
\mathcal{S}^{(h,r)}_{ab}
=
\left|
\frac{\Delta s^{(h,r)}_{ab}}
{\Delta z^{(h)}_{ab}}
\right|
=
\left|
\hat k^{(r)}_a x_b
\alpha^{(h)}_{a,\gamma(b)}
\right|.
\end{equation}
This quantity combines the participation of key coordinate $a$ in the fixed shared head, the representative activation coordinate $b$, and the actual dequantized distance associated with a one-level query-integer move. Including the group scale is essential: $|\hat k^{(r)}_a x_b|$ alone is not the sensitivity with respect to the integer parameter.

Let $\mathcal{B}^{(h)}_a$ contain the input columns whose query weights belong to nondegenerate quantization groups and admit at least one legal one-level move. Frozen or otherwise ineligible coordinates are assigned zero sensitivity and excluded from candidate-region construction. Integer-level sensitivities are then aggregated over eligible input columns with an $\ell_2$ norm:
\begin{align}
\mathcal{S}^{(h,r)}_a
&=
\left(
\sum_{b\in\mathcal{B}^{(h)}_a}
\left(\mathcal{S}^{(h,r)}_{ab}\right)^2
\right)^{1/2} \\
&=
\left|\hat k^{(r)}_a\right|
\left(
\sum_{b\in\mathcal{B}^{(h)}_a}
\left(\alpha^{(h)}_{a,\gamma(b)}\right)^2x_b^2
\right)^{1/2}.
\end{align}
This dimension-level quantity is an $\ell_2$-aggregated sensitivity heuristic used to choose where to search. It is not the exact benefit of a legal move and is not an optimal correction-capacity measure. Every admitted integer candidate must still satisfy the exact loss-decrease condition.

\subsection{Kneedle-Based Protection and Candidate Region}

For each pair $(h,r)$, let $\pi_{h,r}(t)$ be the original query-output index at sorted position $t$, with sensitivities ordered decreasingly:
\begin{equation}
\mathcal{S}^{(h,r)}_{\pi_{h,r}(1)}
\ge
\mathcal{S}^{(h,r)}_{\pi_{h,r}(2)}
\ge\cdots\ge
\mathcal{S}^{(h,r)}_{\pi_{h,r}(d_h)}.
\end{equation}
Before applying Kneedle \cite{satopaa2011kneedle}, both axes are normalized to the unit interval. For $d_h>1$, define normalized rank
\begin{equation}
u_t=\frac{t-1}{d_h-1},
\qquad t=1,\ldots,d_h,
\end{equation}
and, when the sensitivity range is nonzero, normalized decreasing sensitivity
\begin{equation}
y_t
=
\frac{%
\mathcal{S}^{(h,r)}_{\pi_{h,r}(t)}
-
\mathcal{S}^{(h,r)}_{\pi_{h,r}(d_h)}
}{%
\mathcal{S}^{(h,r)}_{\pi_{h,r}(1)}
-
\mathcal{S}^{(h,r)}_{\pi_{h,r}(d_h)}
}.
\end{equation}
Kneedle is configured for a decreasing, convex curve $(u_t,y_t)$ with sensitivity parameter $\kappa$. Let $\rho_{h,r}$ denote the detected knee position. The knee itself is included in the protected extreme region:
\begin{equation}
\mathcal{P}_{h,r}
=
\{
\pi_{h,r}(1),\ldots,\pi_{h,r}(\rho_{h,r})
\}.
\end{equation}
With moderate-region width $H_R$, the allowed query-output rows are
\begin{equation}
\mathcal{M}_{h,r}
=
\{
\pi_{h,r}(\rho_{h,r}+1),\ldots,
\pi_{h,r}(\min(\rho_{h,r}+H_R,d_h))
\}.
\end{equation}

This definition makes the normalization, curve direction, and treatment of the knee explicit. If $d_h\leq1$, all dimension sensitivities are equal, the normalized sensitivity denominator is zero, or Kneedle reports no knee under the selected $\kappa$, the conservative fallback is to skip ReFlip for query head $h$. The fallback avoids inventing an arbitrary moderate region when the sensitivity curve supplies no defensible transition.

Kneedle determines only a candidate region. It does not certify that the protected set, moderate set, or final integer configuration is optimal. Within $\mathcal{M}_{h,r}$, integer-level sensitivities may order candidates, but the exact $\Delta\mathcal{L}^{(h,r)}_{ab}$ remains the sole acceptance criterion.

\subsection{Greedy ReFlip Procedure}

ReFlip proceeds query head by query head. For $h\in\mathcal{G}_r$, it begins from current $\hat W_Q^{(h)}$, fixed $\hat W_K^{(r)}$, and exact residual $e_{h,r}$. It computes integer-level and dimension-level sensitivities, applies the normalized decreasing-curve Kneedle rule, protects $\mathcal{P}_{h,r}$, and searches only $\mathcal{M}_{h,r}$ when a valid knee exists.

For each legal candidate $(a,b,\Delta z^{(h)}_{ab})$, the method evaluates
\begin{equation}
\Delta s^{(h,r)}_{ab}
=
\hat k^{(r)}_a x_b
\alpha^{(h)}_{a,\gamma(b)}
\Delta z^{(h)}_{ab}
\end{equation}
and
\begin{equation}
\Delta\mathcal{L}^{(h,r)}_{ab}
=
e_{h,r}\Delta s^{(h,r)}_{ab}
+\frac{1}{2}\left(\Delta s^{(h,r)}_{ab}\right)^2.
\end{equation}
Candidates with nonnegative exact loss change are rejected. When both symmetric directions are legal and $\mathcal{S}^{(h,r)}_{ab}>0$, at most one can strictly reduce the quadratic loss because the sum of the two exact loss changes is $\left(\mathcal{S}^{(h,r)}_{ab}\right)^2>0$. The unique improving direction, when it exists, is accepted and all affected current-state quantities are updated before another candidate is evaluated. The procedure stops at its accepted-change budget, its pass budget, or a pass with no accepted candidate.

The exact acceptance condition has the equivalent form
\begin{equation}
\left|e_{h,r}+\Delta s^{(h,r)}_{ab}\right|
<
|e_{h,r}|.
\end{equation}
Thus, ReFlip accepts a change only when it moves the current per-head surrogate closer to its full-precision scalar target. The quadratic term rejects a step that overshoots the target by more than it corrects.

\subsection{Scope and Limitations of ReFlip}

ReFlip directly targets a query--key interaction that Flip does not model. A residual caused partly by $\hat W_K^{(r)}-W_K^{(r)}$ can be reduced by a compensating change in $\hat W_Q^{(h)}$. This is a functional correction: the selected query weight need not become closer to its full-precision value. In particular, an accepted ReFlip change may increase the query-head projection residual
\begin{equation}
J_Q^{(h)}
=
\left\|
x\left(\hat W_Q^{(h)}-W_Q^{(h)}\right)^\top
\right\|_2^2
\end{equation}
while strictly reducing $\mathcal{L}_{h,r}$. The two objectives measure different quantities, so ReFlip makes no monotonicity claim for $J_Q^{(h)}$.

The candidate surrogate change and loss change are exact only for the defined unscaled representative-vector scalar. ReFlip does not directly preserve scores between distinct token rows, scaled logits, the softmax distribution, or the value-weighted attention output. A lower per-head scalar loss is evidence only for that targeted criterion.

ReFlip is also greedy and discrete. Candidate ordering and the Kneedle-selected region can affect the endpoint, and a sequence of individually rejected changes may be beneficial if applied jointly. Kneedle does not remove this limitation and does not provide an optimality certificate. ReFlip therefore guarantees strict local improvement only for every accepted change under the stated scalar objective; it does not guarantee a globally optimal query matrix.

\section{Unified Refinement Procedure}

The complete methodology follows a consistent sequence from initialization to local projection correction and then to scalar interaction correction.

\subsection{Stage 1: Calibration and Quantization}

First, calibration activations are summarized into the empirical statistic $\tilde{x}$ and then shrunk toward target $m$ to obtain the representative activation $x$. Next, each selected projection matrix is partitioned along the input dimension into quantization groups using the mapping $\gamma$. Group-specific asymmetric INT4 scales and zero points are computed from zero-inclusive ranges, and RTN initializes the stored integers. This stage produces dequantized query-, key-, and value-head approximations under a common legal integer representation while preserving the distinct GQA head-sharing structure.

\subsection{Stage 2: Flip}

For each projection selected for Flip, compute
\begin{equation}
E=\hat W-W,\qquad r=xE^\top,\qquad J=\|r\|_2^2.
\end{equation}
Evaluate legal one-step integer candidates. For a candidate at $(j,i)$, use
\begin{equation}
\Delta J=2r_jx_i\delta+(x_i\delta)^2,
\qquad
\delta=\alpha_{j,\gamma(i)}\Delta z.
\end{equation}
Accept only a strictly negative exact objective change and immediately update $r_j$. Repeat until the Flip stopping rule is met. The result is a projection matrix whose aggregate residual under $x$ is no larger than at the beginning of the accepted sequence.

\subsection{Stage 3: ReFlip}

After the chosen Flip refinements, process every query head $h\in\mathcal{G}_r$ against its fixed shared key head $r$. Compute
\begin{equation}
s_{h,r}^\star
=
\left[x(W_Q^{(h)})^\top\right]\cdot
\left[x(W_K^{(r)})^\top\right],
\qquad
\hat{s}_{h,r}
=
\left[x(\hat W_Q^{(h)})^\top\right]\cdot
\left[x(\hat W_K^{(r)})^\top\right],
\end{equation}
and set $e_{h,r}=\hat{s}_{h,r}-s_{h,r}^\star$. Hold $\hat W_K^{(r)}$ fixed for every query head in $\mathcal{G}_r$. Compute integer-level and dimension-level sensitivities from current $\hat k^{(r)}$, apply the normalized decreasing-curve Kneedle configuration, include the detected knee in the protected region, and admit the moderate region immediately after it. If the curve has no valid knee, skip ReFlip for that query head.

For every allowed query candidate, evaluate
\begin{equation}
\Delta s^{(h,r)}_{ab}
=
\hat k^{(r)}_a x_b
\alpha^{(h)}_{a,\gamma(b)}
\Delta z^{(h)}_{ab}
\end{equation}
and
\begin{equation}
\Delta\mathcal{L}^{(h,r)}_{ab}
=
e_{h,r}\Delta s^{(h,r)}_{ab}
+\frac{1}{2}\left(\Delta s^{(h,r)}_{ab}\right)^2.
\end{equation}
Accept only $\Delta\mathcal{L}^{(h,r)}_{ab}<0$, update the per-head scalar residual after each accepted change, and terminate under the ReFlip stopping rule. No key or value integer is modified during this stage.

\subsection{Monotonicity of Accepted Changes}

The two stages have simple monotonicity properties under their own stated objectives. Every accepted Flip change satisfies
\begin{equation}
J_{\mathrm{after}}<J_{\mathrm{before}},
\end{equation}
and every accepted ReFlip change satisfies
\begin{equation}
\mathcal{L}^{(h,r)}_{\mathrm{after}}
<
\mathcal{L}^{(h,r)}_{\mathrm{before}}.
\end{equation}
These properties follow from exact objective differences and immediate residual updates. They do not imply monotonic improvement in a different metric, such as unweighted weight reconstruction error, the complete attention matrix, language-model loss, or downstream task quality.

The order of the stages matters conceptually. Flip first improves each selected projection under an activation-weighted output objective. ReFlip then observes each per-query-head surrogate residual left by the current quantized query and shared key heads and uses query-only compensation. Reversing the stages would change the scalar residual and sensitivity profile against which ReFlip operates.

\section{Hyperparameters and Design Choices}

The methodology exposes its principal choices as explicit hyperparameters.

\begin{itemize}
\item \textbf{Bit width $B$.} The present formulation fixes $B=4$, giving the legal unsigned integer set $\mathcal{Z}_4=\{0,\ldots,15\}$. The symbol $B$ is distinct from input-column index $b$.

\item \textbf{Quantization-group size $G$.} The group size determines the column-to-group mapping $\gamma(b)$ and the granularity of scales and zero points. Smaller groups increase quantization metadata but can better adapt to local weight ranges. This quantity is distinct from the GQA head set $\mathcal{G}_r$.

\item \textbf{Calibration sample count $N_{\mathrm{cal}}$.} This controls the amount of activation data used to estimate $\tilde{x}$, the shrinkage target, and the noise level.

\item \textbf{James--Stein target $m$ and noise estimate $\hat{\sigma}^2$.} These determine the shrinkage factor $\lambda$ and therefore the representative activation used by both exact objectives.

\item \textbf{Flip pass budget $T_{\mathrm{F}}$.} This is the maximum number of scans over eligible Flip coordinates.

\item \textbf{Flip acceptance budget $B_{\mathrm{F}}$.} This limits the total number of accepted integer changes for a projection or layer.

\item \textbf{Flip improvement tolerance $\varepsilon_{\mathrm{F}}$.} In finite precision, the strict mathematical condition $\Delta J<0$ may be implemented conservatively as $\Delta J<-\varepsilon_{\mathrm{F}}$.

\item \textbf{ReFlip Kneedle sensitivity $\kappa$.} This controls how the knee detector responds to curvature in the sorted dimension-sensitivity curve.

\item \textbf{Moderate-region width $H_R$.} This determines how many sorted dimensions immediately after the knee are eligible for ReFlip.

\item \textbf{ReFlip pass budget $T_{\mathrm{R}}$ and acceptance budget $B_{\mathrm{R}}$.} These bound the greedy search over legal query-integer changes.

\item \textbf{ReFlip improvement tolerance $\varepsilon_{\mathrm{R}}$.} A conservative finite-precision rule accepts only $\Delta\mathcal{L}^{(h,r)}_{ab}<-\varepsilon_{\mathrm{R}}$.

\item \textbf{Candidate order.} Flip candidates may be ordered by potential activation-weighted step magnitude, while ReFlip candidates may be ordered by $\mathcal{S}^{(h,r)}_{ab}$. Ordering affects the greedy path but never replaces the exact acceptance tests.
\end{itemize}

The one-step neighborhood $\Delta z\in\{-1,+1\}$ is part of the method rather than an unconstrained hyperparameter. It keeps changes local on the existing INT4 grid and makes both objective differences exact and inexpensive to evaluate. Multiple accepted one-step changes may eventually move an integer by more than one level if the candidate is revisited and every intermediate change remains legal and strictly beneficial.

Hyperparameter selection should respect the distinction between refinement opportunity and refinement risk. Very large budgets or a wide post-knee region create more opportunities to reduce the calibration objective but can increase dependence on the representative activation. Conservative budgets, shrinkage, protected extreme dimensions, and positive improvement tolerances are therefore used as design measures intended to limit such over-specialization; their protective effect is not presented as a formal guarantee.

\section{Theoretical Discussion}

\subsection{From Scalar Reconstruction to Functional Residuals}

The methodology is motivated by a hierarchy of objectives. RTN minimizes isolated scalar weight error for fixed quantization parameters. Flip moves to a projection-level functional residual,
\begin{equation}
x(\hat W-W)^\top,
\end{equation}
which accounts for activation magnitude, error sign, and cancellation within each output row. For query head $h\in\mathcal{G}_r$, ReFlip then moves to the per-head scalar surrogate objective
\begin{equation}
\mathcal{L}_{h,r}
=
\frac{1}{2}
\left[
\hat q^{(h)}\cdot\hat k^{(r)}
-
q^{(h)}\cdot k^{(r)}
\right]^2,
\end{equation}
which accounts for the coupling between one query head and its shared GQA key head while constraining changes to the quantized query-head matrix.

Each level deliberately sacrifices breadth for tractability. The objectives are compact enough to permit exact evaluation of every legal one-step candidate. This avoids the ambiguity of claiming that a local scalar change optimizes a large downstream object through an unverified approximation. The trade-off is that exactness applies only to the specified representative objectives.

\subsection{Why Query-Only Compensation Can Help}

The surrogate residual can be expanded conceptually as the difference between a dot product of two quantized projected vectors and the corresponding full-precision dot product. Its value reflects errors in both vectors and their interaction. Once $\hat W_K^{(r)}$ is fixed, changing $\hat W_Q^{(h)}$ alters one side of the dot product while leaving the other side as the reference direction for correction. A suitable query-head change can therefore offset part of the surrogate error associated with the fixed shared quantized key head.

This compensation does not require the modified query weight to approach its original value. It requires only that the exact change in $\hat{s}_{h,r}$ reduce $|\hat{s}_{h,r}-s_{h,r}^\star|$. ReFlip is thus aligned with the same general principle as Flip: a functional aggregate can improve even when an individual scalar reconstruction becomes worse. The query projection residual can likewise become worse because it is not the ReFlip acceptance objective.

\subsection{Role of Sensitivity Protection}

Sensitivity motivates two competing design considerations. A high-sensitivity parameter can correct a residual with a small number of integer changes, but its larger score step can also overshoot the current scalar target. ReFlip therefore does not simply choose the largest sensitivities. As a conservative design choice intended to reduce dependence on extreme score steps and on the representative calibration vector, it protects dimensions in the extreme pre-knee region and searches the moderate region immediately after the knee.

This policy is intentionally risk-aware rather than optimal. The Kneedle partition adapts to the shape of each query head's dimension-sensitivity distribution, while the exact loss test ensures that an admitted candidate is accepted only when it improves the current scalar objective. The separation between heuristic admission and exact acceptance is central: sensitivity controls where to search, whereas $\Delta\mathcal{L}^{(h,r)}_{ab}$ determines whether to change an integer.

\subsection{Boundaries of the Claims}

Several boundaries follow directly from the formulation. First, the representative activation is a James--Stein-shrunk calibration summary, not the complete activation distribution. Second, Flip exactly reduces only its projection residual objective for accepted changes. Third, ReFlip exactly reduces only its per-query-head representative-vector Q--K surrogate loss for accepted changes. Fourth, ReFlip modifies $\hat W_Q^{(h)}$ only and makes no update to $\hat W_K^{(r)}$ or $\hat W_V^{(r)}$.

Most importantly, the ReFlip scalar surrogate does not preserve the full pairwise attention matrix. In a complete sequence, queries and keys arise from different token rows, and attention computation includes many cross-token dot products. A single unscaled surrogate based on $x$ cannot constrain all of them. The method also does not directly optimize the softmax transformation or the value-weighted attention output. These limitations do not invalidate the targeted objectives, but they prevent stronger claims than the mathematics supports.

\subsection{Discrete Search and Optimality}

Both Flip and ReFlip perform greedy coordinate search on a finite integer grid. Their accepted changes are monotonic under their respective objectives, and termination is guaranteed under finite budgets. Without an explicit budget, strict decrease on a finite state space also rules out cycles, although repeated scans may be undesirable.

Neither procedure guarantees a globally optimal quantized matrix. Flip can miss coordinated changes that are beneficial only when applied together. ReFlip can miss a sequence whose first step is neutral or unfavorable but whose combined effect is favorable. Candidate ordering, budget choices, and the Kneedle-selected region can all influence the endpoint. Kneedle itself is a heuristic detector, not an optimality guarantee. The methodological claim is therefore limited to exact evaluation and strict improvement of accepted local changes under the stated objectives.

\section{Summary}

This chapter has defined a post-training refinement methodology for group-wise asymmetric INT4 attention projections under a consistent row-vector convention. A representative James--Stein-shrunk activation $x$ provides the calibration signal. The initial RTN quantizer assigns legal integers $z_{ab}$ with quantization-group-specific scales $\alpha_{a,\gamma(b)}$ and zero points $\zeta_{a,\gamma(b)}$. Zero-inclusive ranges define one-sided groups without ambiguity, while all-zero degenerate groups use an explicit exact frozen representation.

Flip refines a projection matrix through the exact objective
\begin{equation}
J=\left\|x(\hat W-W)^\top\right\|_2^2.
\end{equation}
For a legal dequantized change $\delta$ at $(j,i)$, it evaluates
\begin{equation}
\Delta J=2r_jx_i\delta+(x_i\delta)^2
\end{equation}
and accepts only strict reductions. RTN remains the minimizer of isolated scalar weight error; Flip can improve the aggregate activation-weighted residual by allowing signed error cancellation, even when an individual weight becomes less accurate.

For each query head $h\in\mathcal{G}_r$, ReFlip holds shared $\hat W_K^{(r)}$ fixed and updates $\hat W_Q^{(h)}$ only. It targets the unscaled representative-vector Q--K surrogate loss
\begin{equation}
\mathcal{L}_{h,r}
=\frac{1}{2}
\left[
\hat q^{(h)}\cdot\hat k^{(r)}
-
q^{(h)}\cdot k^{(r)}
\right]^2.
\end{equation}
A legal unit integer change has exact surrogate effect
\begin{equation}
\Delta s^{(h,r)}_{ab}
=
\hat k^{(r)}_a x_b
\alpha^{(h)}_{a,\gamma(b)}
\Delta z^{(h)}_{ab},
\end{equation}
integer-level sensitivity
\begin{equation}
\mathcal{S}^{(h,r)}_{ab}
=
\left|
\hat k^{(r)}_a x_b\alpha^{(h)}_{a,\gamma(b)}
\right|,
\end{equation}
and dimension sensitivity
\begin{equation}
\mathcal{S}^{(h,r)}_a
=
\left|\hat k^{(r)}_a\right|
\left(
\sum_{b\in\mathcal{B}^{(h)}_a}
\left(\alpha^{(h)}_{a,\gamma(b)}\right)^2x_b^2
\right)^{1/2}.
\end{equation}
Kneedle is applied per query head to a unit-normalized decreasing, convex sensitivity curve. It includes the detected knee in the protected extreme region, admits a moderate region immediately after it, and conservatively skips a head when no valid knee exists. Each admitted candidate is evaluated using the exact loss change
\begin{equation}
\Delta\mathcal{L}^{(h,r)}_{ab}
=
e_{h,r}\Delta s^{(h,r)}_{ab}
+\frac{1}{2}\left(\Delta s^{(h,r)}_{ab}\right)^2,
\end{equation}
and the per-head scalar residual is updated after every accepted change. Such a change may worsen query projection residual while reducing the targeted scalar loss.

{\raggedright
The resulting procedure retains the theoretical memory advan\-tages and deploy\-ment form of group-wise INT4 weight-only quantization. Its guarantees are deliberately limited to the stated objectives. Kneedle is heuristic, greedy refinement is not globally optimal, and the scalar ReFlip surrogate does not preserve the full pairwise attention matrix. Within these boundaries, Flip first corrects the projection residual under the representative activation, and ReFlip then performs exact per-query-head compensation of the unscaled representative-vector Q--K surrogate.
\par}
